\documentclass[11pt,a4paper]{article}
\usepackage[utf8]{inputenc}
\usepackage[T1]{fontenc}
\usepackage{lmodern}
\usepackage[french]{babel}
\usepackage{amsmath,amssymb,mathtools}
\usepackage{icomma}
\usepackage{xcolor}
\usepackage{colortbl}
\usepackage{tikz}
\usepackage[most]{tcolorbox}
\usepackage{enumitem}
\usepackage{array}
\usepackage{booktabs}
\usepackage{fancyhdr}
\usepackage[hidelinks]{hyperref}
\usetikzlibrary{arrows.meta,positioning,calc,patterns,backgrounds,shapes.geometric,decorations.pathreplacing,decorations.markings}
\usepackage[a4paper,margin=2.15cm,headheight=26pt,headsep=12pt]{geometry}
\usepackage{titlesec}

% ---------------- Palette ----------------
\definecolor{primary}{RGB}{0,151,169}
\definecolor{primarydark}{RGB}{0,88,101}
\definecolor{defcol}{RGB}{0,114,178}
\definecolor{thmcol}{RGB}{0,140,120}
\definecolor{formulacol}{RGB}{198,93,9}
\definecolor{excol}{RGB}{0,135,90}
\definecolor{attncol}{RGB}{200,40,55}
\definecolor{methcol}{RGB}{90,90,100}
\definecolor{remarkcol}{RGB}{120,94,200}
\definecolor{barcol}{RGB}{0,151,169}
\definecolor{barcold}{RGB}{0,110,124}
\definecolor{modecol}{RGB}{200,55,55}
\definecolor{meancol}{RGB}{0,90,200}
\definecolor{medcol}{RGB}{160,90,190}
\definecolor{quartcol}{RGB}{0,135,90}
\definecolor{ogivecol}{RGB}{176,74,0}

% ---------------- Boîtes ----------------
\tcbset{boxbase/.style={enhanced,breakable,boxrule=0.9pt,arc=2.5pt,
  left=3mm,right=3mm,top=2.3mm,bottom=2.3mm,
  fonttitle=\bfseries,coltitle=white,toptitle=0.6mm,bottomtitle=0.6mm}}

\newtcolorbox{definition}[1][]{boxbase,colback=defcol!6,colframe=defcol,
  title={\textsc{Définition}\ifx\relax#1\relax\relax\else\textmd{ : \itshape #1}\fi}}
\newtcolorbox{propriete}[1][]{boxbase,colback=thmcol!7,colframe=thmcol,
  title={\textsc{Propriété}\ifx\relax#1\relax\relax\else\textmd{ : \itshape #1}\fi}}
\newtcolorbox{formule}[1][]{boxbase,colback=formulacol!7,colframe=formulacol,
  title={\textsc{Formule}\ifx\relax#1\relax\relax\else\textmd{ : \itshape #1}\fi}}
\newtcolorbox{exemple}[1][]{boxbase,colback=excol!6,colframe=excol,
  title={\textsc{Exemple}\ifx\relax#1\relax\relax\else\textmd{ : \itshape #1}\fi}}
\newtcolorbox{methode}[1][]{boxbase,colback=methcol!8,colframe=methcol,sharp corners,
  title={\textsc{Méthode}\ifx\relax#1\relax\relax\else\textmd{ : \itshape #1}\fi}}
\newtcolorbox{remarque}[1][]{boxbase,colback=remarkcol!7,colframe=remarkcol,
  title={\textsc{Astuce}\ifx\relax#1\relax\relax\else\textmd{ : \itshape #1}\fi}}
\newtcolorbox{attention}[1][]{boxbase,colback=attncol!7,colframe=attncol,
  title={\textsc{Attention}\ifx\relax#1\relax\relax\else\textmd{ : \itshape #1}\fi}}

% Boîte "exercice" et "corrigé" (titre libre passé en argument)
\newtcolorbox{exobox}[1]{boxbase,colback=primary!4,colframe=primary,
  before skip=8pt,after skip=8pt,title={#1}}
\newtcolorbox{corbox}[1]{boxbase,colback=excol!5,colframe=excol!85!black,
  before skip=8pt,after skip=8pt,title={#1}}

% ---------------- En-tête ----------------
\newcommand{\docpart}[1]{\def\thedocpart{#1}}
\docpart{}
\pagestyle{fancy}
\fancyhf{}
\fancyhead[L]{\small\textcolor{primarydark}{\textbf{Fiche 8} — Statistiques descriptives}}
\fancyhead[R]{\small\textcolor{primarydark}{\textsc{\thedocpart}}}
\fancyfoot[C]{\small\thepage}
\renewcommand{\headrulewidth}{0.8pt}
\renewcommand{\headrule}{\hbox to\headwidth{\color{primary}\leaders\hrule height \headrulewidth\hfill}}

\titleformat{\section}{\large\bfseries\color{primarydark}}{\thesection}{0.6em}{}
\titleformat{\subsection}{\normalsize\bfseries\color{primary}}{\thesubsection}{0.6em}{}
\setlist{topsep=2pt,itemsep=1.5pt,parsep=0pt}

% Bandeau de titre du document
\newcommand{\bandeau}[2]{%
\begin{tcolorbox}[boxbase,colback=primary,colframe=primarydark,halign=center,
   before skip=2pt,after skip=12pt]
{\LARGE\bfseries\color{white} #1}\\[3pt]{\normalsize\color{white} #2}
\end{tcolorbox}}

\newcommand{\ecm}{\sigma_m}
\newcommand{\ect}{\sigma}
\newcommand{\med}{M\!e}
\docpart{Tutoriel — Thème 2}
\begin{document}
\bandeau{Thème 2 — Mode et moyenne}{Valeur la plus fréquente et moyenne pondérée}

\noindent\textbf{But du thème.} Deux indicateurs de \emph{position} (le « centre » d'une série) :
le \textbf{mode}, qui se lit sans calcul, et la \textbf{moyenne}, qui se calcule toujours de la même
façon — une \emph{moyenne pondérée}. Ce thème automatise ces deux gestes et la propriété de
\textbf{linéarité} de la moyenne (très rentable au concours).

\section{Le mode : la valeur la plus fréquente}

\begin{definition}[mode $M$]
Le \textbf{mode} est la modalité (ou la classe) dont l'\textbf{effectif $n_i$ est le plus grand}.
On le lit directement dans la colonne des effectifs : aucun calcul. Attention, le mode est une
\emph{valeur} (ou une classe), pas un effectif.
\end{definition}

\begin{exemple}[lecture du mode]
Notes : $2$ ($15$ élèves), $5$ ($10$ élèves), $8$ ($5$ élèves). Le plus grand effectif est $15$,
donc $M=2$ (et surtout pas $M=15$). Sur des classes, si $[85,95[$ a l'effectif maximal, alors
$M=[85,95[$ (ou $M=90$, son centre).
\end{exemple}

\section{La moyenne arithmétique pondérée}

\begin{formule}[moyenne d'une série regroupée]
\[
\boxed{\;m=\frac{1}{N}\sum_{i=1}^{k} n_i\,x_i=\frac{n_1x_1+n_2x_2+\cdots+n_kx_k}{N}\;}
\]
$x_i$ : centre de classe (ou valeur si discret) ; $n_i$ : effectif (le « poids ») ; $N=\sum n_i$.
\textbf{Unité :} celle du caractère. On multiplie \emph{toujours} chaque valeur par son effectif :
c'est une moyenne \textbf{pondérée}, pas une moyenne « simple » des $x_i$.
\end{formule}

\begin{methode}[calculer une moyenne en 3 gestes]
\begin{enumerate}
  \item Dresser la colonne des produits $n_i x_i$ (valeur $\times$ effectif).
  \item Additionner cette colonne : on obtient $\sum n_i x_i$.
  \item Diviser par $N$.
\end{enumerate}
\end{methode}

\begin{exemple}[moyenne pondérée]
Notes $2$ ($15$), $5$ ($10$), $8$ ($5$), $N=30$ :
\[
m=\frac{15\times2+10\times5+5\times8}{30}=\frac{30+50+40}{30}=\frac{120}{30}=\mathbf{4}.
\]
La note $2$, très fréquente, « tire » la moyenne vers le bas : c'est tout l'effet des poids.
\end{exemple}

\section{La moyenne est linéaire : le réflexe qui fait gagner du temps}

\begin{propriete}[effet d'un changement d'échelle]
Si l'on transforme chaque valeur par $x_i\mapsto \alpha x_i+\beta$, alors la moyenne devient
\[
\boxed{\,m'=\alpha\,m+\beta\,.}
\]
Cas particuliers : ajouter $\beta$ à toutes les valeurs ajoute $\beta$ à la moyenne ;
multiplier toutes les valeurs par $\alpha$ multiplie la moyenne par $\alpha$.
\end{propriete}

\begin{exemple}[atteindre une moyenne cible]
Série de moyenne $m=4/10$. On veut une moyenne de $5$. En ajoutant $1$ point à chaque note
($\alpha=1,\ \beta=1$) : $m'=1\times4+1=5$. \checkmark\;
En \emph{doublant} chaque note ($\alpha=2,\ \beta=0$) on obtiendrait $m'=2\times4=8$.
\end{exemple}

\begin{remarque}[repérer le mode et la moyenne « à la main »]
La moyenne peut tomber \emph{entre} deux valeurs (ex.\ $4{,}0$ n'est la note de personne) : c'est
normal, c'est un résumé, pas une valeur observée. Le mode, lui, est \emph{toujours} une valeur
réellement présente (celle du plus grand bâton / de la plus haute barre).
\end{remarque}

\begin{attention}[ne pas oublier de pondérer]
L'erreur nº1 : faire $\dfrac{2+5+8}{3}=5$ au lieu de la moyenne pondérée $4$. Tant qu'il y a des
effectifs (ou des classes), \textbf{on pondère}. La moyenne « simple » ne vaut que si tous les
effectifs sont égaux à $1$.
\end{attention}
\end{document}
